\documentclass[aspectratio=169]{beamer}
\usetheme{Madrid}
\usepackage[utf8]{inputenc}
\usepackage[T1]{fontenc}
\usepackage{tikz}
\usepackage{amssymb}
\usepackage{graphicx}
\setbeamertemplate{navigation symbols}{}
\setbeamercolor{title}{fg=white,bg=red!45!black}
\setbeamercolor{frametitle}{fg=white,bg=red!45!black}
\setbeamercolor{block title}{bg=red!45!black!20,fg=black}
\setbeamercolor{block body}{bg=red!45!black!8,fg=black}
\setbeamercolor{item}{fg=red!45!black}
\setbeamercolor{itemize item}{fg=red!45!black}
\setbeamertemplate{itemize item}[triangle]

\title{Win/Lose Conditions}
\subtitle{if/else checks whether health has reached 0 or below}
\author{ByteForge --- Grade 8 --- App Builder}
\date{}

\begin{document}

\begin{frame}
    \titlepage
\end{frame}

\begin{frame}{Today's Plan}
    \begin{itemize}
        \item Game Over?
        \item[$\bigstar$] BOSS --- Game Over?
    \end{itemize}
\end{frame}

\begin{frame}{\texttt{<=} catches EXACTLY 0, not just below it}
    \large
    \begin{itemize}
        \item \texttt{if health <= 0:}
        \item Using \texttt{<} instead would wrongly treat health of exactly 0 as still "alive", a genuine off-by-one bug in a real battle system.
    \end{itemize}
\end{frame}


\begin{frame}{Every branch needs its own outcome}
    \large
    \begin{itemize}
        \item The \texttt{else} branch needs its own \texttt{print()} too
        \item One message for losing, a different one for still fighting, never left blank.
    \end{itemize}
\end{frame}

\begin{frame}{Real-World Examples}
    \large
    \begin{itemize}
        \item Every real game's "you lose" check is exactly this
        \item An inclusive \texttt{<= 0} comparison, because reaching exactly zero health should genuinely count as defeated, not a special edge case that slips through.
    \end{itemize}
\end{frame}


\begin{frame}{Let Them Discover}
    \Large
    If health reaches EXACTLY 0, should that count as a loss? Most people say yes immediately — but \texttt{health < 0} would actually say no. Which operator does the job correctly?
\end{frame}

\begin{frame}{Your Turn}
    \Large
    Open the student map, find \textbf{Win/Lose Conditions}, and work through it
    hands-on --- everything you just saw, they now do themselves.
\end{frame}

\end{document}
